\part{Markets, Welfare, and Market Structures}

\chapter{Competitive Markets and Welfare}
\label{chap:competition}

\section{Learning objectives}

After completing this chapter, you should be able to:
\begin{itemize}
  \item determine competitive market equilibrium;
  \item use comparative statics to analyse shifts in supply and demand;
  \item calculate consumer surplus, producer surplus, and total surplus; and
  \item evaluate the effects of taxes, price ceilings, and price floors.
\end{itemize}

\section{Competitive markets}

A \keyterm{competitive market} has many buyers and sellers, a homogeneous product, and no individual participant large enough to affect market price. Buyers and sellers are price takers. Competition is an ideal type: real markets may approximate it to different degrees.

Market demand gives the quantity consumers wish to buy at each price; market supply gives the quantity firms wish to sell at each price. A \keyterm{market equilibrium} is a price and quantity at which quantity demanded equals quantity supplied:
\[
Q_D(p^*)=Q_S(p^*).
\]
At prices above equilibrium, there is excess supply, placing downward pressure on price. At prices below equilibrium, there is excess demand, placing upward pressure on price. This adjustment story describes a tendency, not necessarily an instantaneous process.

\section{Comparative statics}

\keyterm{Comparative statics} compares equilibria before and after a change in an underlying condition. An increase in demand shifts the demand curve rightward, raising equilibrium price and quantity when supply slopes upward. An increase in production cost shifts supply leftward, raising price and reducing quantity.

Comparative statics is often more valuable than a single equilibrium calculation. It identifies which changes are expected to affect prices, quantities, and distribution, while leaving open the empirical question of their magnitude.

\section{Consumer surplus and producer surplus}

\keyterm{Consumer surplus} is the difference between consumers' willingness to pay and the amount they actually pay. In a demand diagram, it is the area below demand and above market price for units traded. \keyterm{Producer surplus} is the difference between the price received and the minimum amount required for firms to supply each unit; it is the area above supply and below market price.

For linear curves, these areas are often triangles. If demand intercepts the price axis at $P_{max}$, equilibrium price is $p^*$, and equilibrium quantity is $Q^*$, consumer surplus is
\[
CS=\frac{1}{2}(P_{max}-p^*)Q^*.
\]

\section{Efficiency and total surplus}

\keyterm{Total surplus} is consumer surplus plus producer surplus. In a competitive market without external effects, equilibrium maximises total surplus: units for which willingness to pay exceeds marginal cost are traded, while units for which marginal cost exceeds willingness to pay are not.

This conclusion is conditional. It requires that prices reflect relevant costs and benefits, that market participants have suitable information, and that the distribution of purchasing power is not itself the criterion for justice. Later chapters study cases in which these conditions fail.

\section{Taxes}

A per-unit tax creates a wedge between the price paid by buyers and the price received by sellers. The quantity traded falls, tax revenue is collected by government, and part of the potential surplus from mutually beneficial trade disappears. This loss is \keyterm{deadweight loss}.

The legal incidence of a tax does not determine its economic incidence. Whether the tax is formally paid by buyers or sellers, its burden is shared according to relative elasticities. The less elastic side of the market bears more of the burden because it has fewer alternatives.

\section{Price controls}

A \keyterm{price ceiling} is a legal maximum price. If it is set below equilibrium, quantity demanded exceeds quantity supplied, creating a shortage. Allocation may then occur through waiting, queues, informal payments, rationing, or discrimination rather than price.

A \keyterm{price floor} is a legal minimum price. If it is set above equilibrium, quantity supplied exceeds quantity demanded, creating a surplus. Policy design must then address what happens to unsold output or unemployed resources.

Price controls can pursue equity or stability objectives, but their effects cannot be assessed only by observing the controlled price. The quantity available, method of allocation, quality, and incentives to enter or exit the market also matter.

\section*{Chapter Review}

\subsection*{Key concepts}
\begin{itemize}
  \item \textbf{Market equilibrium}: $Q_D=Q_S$.
  \item \textbf{Consumer surplus}: willingness to pay minus expenditure.
  \item \textbf{Producer surplus}: revenue minus variable cost, represented by the area above supply.
  \item \textbf{Total surplus}: consumer plus producer surplus.
  \item \textbf{Deadweight loss}: surplus lost because beneficial trades do not occur.
  \item \textbf{Tax incidence}: the actual division of a tax burden.
\end{itemize}

\section*{Exercises}

\subsection*{Conceptual questions}
\begin{enumerate}[label=\arabic*.]
  \item Define competitive market equilibrium. What creates pressure for price adjustment away from equilibrium?
  \item Distinguish consumer surplus from producer surplus.
  \item Why does a competitive equilibrium maximise total surplus only under particular assumptions?
  \item Explain why legal tax incidence need not equal economic tax incidence.
\end{enumerate}

\subsection*{Graphical exercises}
\begin{enumerate}[resume,label=\arabic*.]
  \item Draw supply and demand. Mark equilibrium, consumer surplus, producer surplus, and total surplus.
  \item Show the effect of a per-unit tax. Mark prices paid and received, tax revenue, and deadweight loss.
  \item Draw a binding price ceiling and identify the shortage.
\end{enumerate}

\subsection*{Analytical exercises}
\begin{enumerate}[resume,label=\arabic*.]
  \item Let $Q_D=100-2P$ and $Q_S=10+P$. Find the competitive equilibrium.
  \item Using Exercise 7, calculate consumer and producer surplus.
  \item A tax of \euro6 per unit is imposed in the market of Exercise 7. Find the new quantity, buyer price, seller price, and tax revenue.
\end{enumerate}

\subsection*{Exam-style multiple-choice questions}
\begin{enumerate}[resume,label=\arabic*.]
  \item A binding price ceiling is set: (a) above equilibrium price; (b) below equilibrium price; (c) exactly at every possible price; or (d) only in monopolies.
  \item A tax generally creates deadweight loss because it: (a) makes all trade impossible; (b) reduces the quantity of mutually beneficial trades; (c) eliminates tax revenue; or (d) raises consumer surplus.
  \item The less elastic side of a market generally bears: (a) less of a tax; (b) more of a tax; (c) no tax; or (d) all producer surplus.
\end{enumerate}

\subsection*{PPE discussion question}

\noindent Competitive equilibrium can be efficient while leaving some people with very little purchasing power. Discuss why efficiency and justice are distinct criteria for evaluating market outcomes.
